% ============================================================
% PI-SSL Paper
% Physics-Informed Self-Supervised Learning for
% Cross-Domain Wi-Fi Gesture Recognition
%
% COMPLETE SECTIONS (backed by verified code output):
%   Section 3 — Dataset and Preprocessing
%   Section 4.1 — Experimental Setup
%
% PLACEHOLDER SECTIONS (filled once results exist):
%   Section 1 — Introduction
%   Section 2 — Related Work
%   Section 4.2+ — Results, Ablation
%   Section 5 — Conclusion
%
% Template: IEEEtran (journal mode)
% Compile:  pdflatex pi_ssl_paper.tex
% ============================================================

\documentclass[journal]{IEEEtran}

% ── Packages ─────────────────────────────────────────────────────────────────
\usepackage{amsmath}
\usepackage{amssymb}
\usepackage{booktabs}
\usepackage{multirow}
\usepackage{graphicx}
\usepackage{xcolor}
\usepackage{url}
\usepackage{cite}
\usepackage{array}
\usepackage{tabularx}
\usepackage{makecell}

% ── Convenience macros ────────────────────────────────────────────────────────
\newcommand{\bvp}{\mathbf{X}}
\newcommand{\R}{\mathbb{R}}

% =============================================================================
\begin{document}
% =============================================================================

\title{Physics-Informed Self-Supervised Learning for\\
Cross-Domain Wi-Fi Gesture Recognition}

\author{Ghulam~Fatima,~\IEEEmembership{Student Member,~IEEE,}
        and~Zahid~Halim,~\IEEEmembership{Member,~IEEE}%
\thanks{G.~Fatima and Z.~Halim are with the Department of Information
Management, National Yunlin University of Science and Technology (YunTech),
Yunlin 640, Taiwan (e-mail: D11423005@yuntech.edu.tw).}}

\maketitle

% =============================================================================
% ABSTRACT  (placeholder — fill once results are known)
% =============================================================================
\begin{abstract}
[To be completed after experiments.]
Wi-Fi-based gesture recognition has emerged as a promising contactless
sensing modality, yet existing methods generalise poorly across physical
environments due to domain shift in channel conditions.
We propose PI-SSL, a physics-informed self-supervised learning framework
that leverages the domain-invariant structure of Body-coordinate Velocity
Profiles (BVP) as a self-supervisory signal.
Under a 3-fold Leave-One-Room-Out evaluation on Widar3.0 using only 25\%
labeled data, PI-SSL achieves [X.X\%] mean accuracy, surpassing the
fully-supervised Widar3.0 TPAMI baseline (89.7\%) that uses 100\% labels.
\end{abstract}

\begin{IEEEkeywords}
Wi-Fi sensing, gesture recognition, self-supervised learning,
cross-domain generalisation, Body-coordinate Velocity Profile.
\end{IEEEkeywords}

% =============================================================================
\section{Introduction}
\label{sec:intro}
% =============================================================================

[To be completed.]

% =============================================================================
\section{Related Work}
\label{sec:related}
% =============================================================================

[To be completed.]

% =============================================================================
\section{Dataset and Preprocessing}
\label{sec:dataset}
% =============================================================================

\subsection{Widar3.0 BVP Dataset}
\label{sec:dataset:overview}

We evaluate on the publicly released Widar3.0
dataset~\cite{zheng2019widar3}, which provides Body-coordinate Velocity
Profile (BVP) signals collected across multiple physical environments
using commodity Wi-Fi hardware.
Each BVP sample is derived from Channel State Information (CSI)
measurements by applying a domain transformation that expresses hand
motion as a velocity profile in a coordinate system fixed to the human
body, partially decoupling the signal from the specific
transmitter--receiver geometry of the recording environment.

The dataset comprises recordings across \textbf{14 session environments}
spanning three distinct physical rooms: a \emph{Classroom}
(Room~1, 5~sessions), a \emph{Hall} (Room~2, 8~sessions), and an
\emph{Office} (Room~3, 1~session).
Seventeen participants (user1--user17) performed gesture sequences across
these sessions.
After filename parsing and gesture label resolution
(Section~\ref{sec:dataset:labels}), we obtain \textbf{43,158 usable
files}; 500 files are discarded because they belong to users whose
gesture mappings are absent from the official dataset documentation for
three partially-recorded environments.

\subsubsection*{File format}
Each \texttt{.mat} file stores one gesture instance as a three-dimensional
array \texttt{velocity\_spectrum\_ro} $\in \R^{20 \times 20 \times T}$,
where the first axis indexes Doppler-velocity bins, the second indexes
spatial/angle bins, and $T \in [0,\,38]$ is the number of recorded time
frames.
The filename encodes per-sample metadata in the format:
\begin{equation*}
  \texttt{\{user\}-\{gesture\}-\{torso\}-\{face\}-\{rep\}-[\ldots]-L\{link\}.mat}
\end{equation*}
where \texttt{gesture} is a local integer label whose semantic meaning
is session-specific (Section~\ref{sec:dataset:labels}).

% ── Gesture label resolution ─────────────────────────────────────────────────
\subsection{Gesture Label Resolution}
\label{sec:dataset:labels}

A critical property of Widar3.0 that is underemphasised in prior work is
that the gesture integers embedded in filenames are \emph{local to each
recording session}: the same integer maps to different semantic classes
across environments.
For example, local label~\texttt{1} encodes \textsc{Push\&Pull} in the
20181109-VS session but \textsc{Draw-1} (digit) in 20181112-VS.
Treating raw filename integers as global class labels introduces spurious
cross-environment label noise and inflates the apparent number of classes.

We resolve every local label to a global semantic name using the complete
per-environment, per-user mapping published in the official Widar3.0
README, yielding \textbf{21 unique semantic gesture classes}.
The digit gesture \textsc{Draw-0} is absent from all recorded sessions.
Table~\ref{tab:gestures} lists the full class inventory and the per-class
sample count after preprocessing.
Only four of the fourteen environments contain all six standard benchmark
gestures; these four environments are the sole contributors to the
supervised training and test sets.

\begin{table}[t]
  \centering
  \caption{Semantic gesture classes in Widar3.0 and sample counts after
           preprocessing. Standard benchmark gestures are marked
           ($\star$).}
  \label{tab:gestures}
  \renewcommand{\arraystretch}{1.15}
  \begin{tabular}{@{}llr@{}}
    \toprule
    \textbf{Group} & \textbf{Gesture} & \textbf{Samples} \\
    \midrule
    \multirow{6}{*}{Benchmark$^\star$}
      & Push\&Pull$^\star$    & 4{,}297 \\
      & Sweep$^\star$         & 4{,}174 \\
      & Clap$^\star$          & 4{,}170 \\
      & Slide$^\star$         & 4{,}875 \\
      & Draw-O-H$^\star$      & 3{,}748 \\
      & Draw-Zigzag-H$^\star$ & 3{,}750 \\
    \midrule
    \multirow{6}{*}{Extended}
      & Draw-N-H         & 3{,}124 \\
      & Draw-Triangle-H  & 2{,}999 \\
      & Draw-Rectangle-H & 3{,}000 \\
      & Draw-O-V         &    675  \\
      & Draw-Zigzag-V    & 1{,}925 \\
      & Draw-N-V         & 1{,}923 \\
    \midrule
    \multirow{9}{*}{Digit}
      & Draw-1 &   500 \\
      & Draw-2 &   500 \\
      & Draw-3 &   500 \\
      & Draw-4 &   500 \\
      & Draw-5 &   500 \\
      & Draw-6 &   492 \\
      & Draw-7 &   500 \\
      & Draw-8 &   500 \\
      & Draw-9 &   500 \\
    \midrule
    \multicolumn{2}{@{}l}{\textbf{Total}} & \textbf{43{,}152} \\
    \bottomrule
  \end{tabular}
\end{table}

% ── Preprocessing ─────────────────────────────────────────────────────────────
\subsection{Preprocessing}
\label{sec:dataset:preprocessing}

\subsubsection*{Per-sample normalisation}
BVP signal amplitude varies substantially across physical environments
due to differences in room geometry, path loss, and receiver placement.
To remove this environment-level offset and encourage the SSL encoder to
learn gesture kinematics rather than absolute signal strength, we apply
per-sample z-score normalisation to every BVP volume before any other
operation:
\begin{equation}
  \hat{\bvp} =
    \frac{\bvp - \mu(\bvp)}{\sigma(\bvp) + \varepsilon},
  \qquad \varepsilon = 10^{-6},
  \label{eq:zscore}
\end{equation}
where $\bvp \in \R^{20 \times 20 \times T}$ is the raw BVP array and
$\mu(\cdot)$, $\sigma(\cdot)$ denote the mean and standard deviation
computed over all elements of the volume.
Normalisation is applied \emph{before} temporal padding so that the
appended zeros retain the semantic of ``no signal'' and do not
participate in the sample statistics.

\subsubsection*{Temporal alignment}
The time dimension $T$ ranges from 0 to 38 across files, reflecting
variation in gesture duration and recording length.
We adopt $T^{*} = 20$ as the canonical sequence length, matching the
mode of the empirical distribution.
Samples with $T > T^{*}$ are truncated to the first $T^{*}$ frames,
retaining the peak of the gesture trajectory; samples with $T < T^{*}$
are zero-padded on the right.
Six files with $T = 0$ or corrupted array dimensionality are excluded,
yielding a final dataset of \textbf{43,152 samples}, each of shape
$20 \times 20 \times 20$.

The preprocessed dataset is stored as a single compressed
\texttt{.npz} archive (159.6\,MB on disk), eliminating per-sample
\texttt{.mat} parsing overhead during GPU training and ensuring full
reproducibility of all splits.

% =============================================================================
\section{Experiments}
\label{sec:experiments}
% =============================================================================

\subsection{Experimental Setup}
\label{sec:experiments:setup}

\subsubsection*{Evaluation protocol}
We follow the Leave-One-Room-Out (LORO) cross-validation protocol
introduced by Zheng~\textit{et al.}~\cite{zheng2019widar3}, which is
the standard benchmark for cross-environment Wi-Fi gesture recognition.
In each of three folds, one physical room is withheld entirely as the
test domain and the remaining two rooms form the training domain.
No environment seen during training appears in evaluation, providing a
stringent measure of cross-domain generalisation.
Evaluation is restricted to the six standard benchmark gestures
(\textsc{Push\&Pull}, \textsc{Sweep}, \textsc{Clap}, \textsc{Slide},
\textsc{Draw-O-H}, \textsc{Draw-Zigzag-H}) to enable direct numerical
comparison with all published results on this dataset.

\subsubsection*{Semi-supervised split}
We simulate a low-annotation setting by using \textbf{25\%} of the
standard-6 training files as the labeled set for supervised fine-tuning.
The split is stratified by gesture class so every class contributes a
proportional number of labeled examples (minimum one per class).
The remaining 75\% of standard-6 training files is treated as unlabeled
and may be used for consistency-based auxiliary losses.
SSL pre-training uses \emph{all} files from the training rooms across
all 21 gesture classes (labels ignored), maximising the volume of Wi-Fi
signal available to the representation encoder.

Table~\ref{tab:splits} reports the exact file counts for each fold.

\begin{table}[t]
  \centering
  \caption{Dataset split sizes per LORO fold. SSL pre-train uses all
           gesture classes; labeled and unlabeled sets contain
           standard-6 gestures only.}
  \label{tab:splits}
  \renewcommand{\arraystretch}{1.15}
  \begin{tabular}{@{}clrrrr@{}}
    \toprule
    \textbf{Fold} & \textbf{Test Room} &
    \textbf{SSL Pre-train} & \textbf{Labeled} &
    \textbf{Unlabeled} & \textbf{Test} \\
    \midrule
    0 & Office    & 40{,}157 & 5{,}505 & 16{,}514 &  2{,}995 \\
    1 & Classroom & 11{,}744 & 2{,}155 &  6{,}464 & 16{,}398 \\
    2 & Hall      & 34{,}410 & 4{,}849 & 14{,}545 &  5{,}623 \\
    \bottomrule
  \end{tabular}
\end{table}

\subsubsection*{Metric}
Performance is reported as mean per-fold classification accuracy
averaged across all three folds, with standard deviation.
The target is ${>}89.5\%$ mean accuracy, exceeding the Widar3.0 TPAMI
baseline of $89.7\%$~\cite{zheng2019widar3} under identical evaluation
conditions while using only 25\% of the available labels.

\subsubsection*{Baselines}
We compare PI-SSL against three published systems:

\begin{itemize}

  \item \textbf{Widar3.0}~\cite{zheng2019widar3}: BVP features with an
  SVM classifier trained on 100\% labeled data under the identical LORO
  protocol and standard-6 gesture set.
  Reported mean accuracy: \textbf{89.7\%}.
  This is the primary baseline; PI-SSL must exceed it using only 25\%
  labels.

  \item \textbf{SenseFi (CNN-5)}~\cite{ma2023sensefi}: A five-layer CNN
  trained on random 80/20 subject-independent splits covering all 22
  gesture classes.
  Reported accuracy: \textbf{70.19\%}.
  Note that this protocol does not enforce cross-environment evaluation;
  the comparison is included for context only and does not constitute a
  fair direct comparison.

  \item \textbf{ARC-Fi}~\cite{yang2023arcfi}: A semi-supervised
  cross-room method that reports 79--91\% accuracy across different
  configurations on Widar3.0.
  Direct comparison is approximate due to differences in gesture set
  size and split strategy.

\end{itemize}

\subsection{Results}
\label{sec:experiments:results}

Table~\ref{tab:results} reports per-fold and mean classification accuracy
for all baselines under the 3-fold LORO protocol with 25\% labeled data.
Rows marked with $\dagger$ show the same method trained on 100\% labels,
providing an internal upper bound for each architecture.

\begin{table}[t]
  \centering
  \caption{Classification accuracy (\%) on Widar3.0 standard-6 gestures,
           3-fold LORO. \textbf{PI-SSL uses only 25\% of labels};
           $\dagger$ denotes the same method with 100\% labels for
           reference. Best 25\%-label result in \textbf{bold}.}
  \label{tab:results}
  \renewcommand{\arraystretch}{1.15}
  \begin{tabular}{@{}lcccc@{}}
    \toprule
    \textbf{Method} &
    \makecell{\textbf{Fold 0}\\\textit{Office}} &
    \makecell{\textbf{Fold 1}\\\textit{Classroom}} &
    \makecell{\textbf{Fold 2}\\\textit{Hall}} &
    \textbf{Mean $\pm$ Std} \\
    \midrule
    \multicolumn{5}{@{}l}{\textit{Supervised upper bounds (100\% labels)}} \\
    BVP+SVM$^\dagger$           & 77.36 & 58.69 & 71.05 & 69.03 $\pm$ 7.76 \\
    Widar\_LeNet$^\dagger$~\cite{ma2023sensefi}    & 85.58 & 68.74 & 83.14 & 79.15 $\pm$ 7.43 \\
    Widar\_CNN-GRU$^\dagger$~\cite{ma2023sensefi}  & 79.87 & 63.84 & 78.20 & 73.97 $\pm$ 7.19 \\
    Widar3.0 TPAMI~\cite{zheng2019widar3}$^*$      & \multicolumn{3}{c}{---} & 89.70 \\
    \midrule
    \multicolumn{5}{@{}l}{\textit{Baselines (25\% labels)}} \\
    BVP+SVM (ours)              & 73.29 & 51.62 & 62.87 & 62.59 $\pm$ 8.85 \\
    Widar\_LeNet~\cite{ma2023sensefi}    & 75.46 & 57.58 & 74.39 & 69.14 $\pm$ 8.19 \\
    Widar\_CNN-GRU~\cite{ma2023sensefi}  & 74.72 & 53.67 & 70.46 & 66.28 $\pm$ 9.09 \\
    SimCLR + fine-tune          & -- & -- & -- & -- \\
    \midrule
    \textbf{PI-SSL (ours)}      & -- & -- & -- & \textbf{--} \\
    \bottomrule
    \multicolumn{5}{@{}l}{\footnotesize $^*$Different feature extraction
      and evaluation setup; not directly comparable (see text).} \\
  \end{tabular}
\end{table}

\subsubsection*{SVM baseline}
The BVP+SVM baseline achieves a mean accuracy of $62.59 \pm 8.85\%$
with 25\% labeled data and $69.03 \pm 7.76\%$ with 100\% labels.
The large inter-fold variance reflects a structural asymmetry in the
LORO protocol: Fold~1 (Classroom) provides only 2{,}155 labeled training
samples against 16{,}396 test samples, while Fold~0 (Office) reverses
this ratio (5{,}505 training, 2{,}995 test), yielding an 21.7-point
accuracy spread across folds.
The Classroom fold consistently produces the lowest accuracy across all
methods, making it the dominant contributor to variance.

We note that our SVM operates on the raw flattened BVP volume
($20 \times 20 \times 20 \to 8{,}000$-dim) with an RBF kernel
($C = 10$, $\gamma = \texttt{scale}$), whereas the Widar3.0 TPAMI
paper~\cite{zheng2019widar3} reports 89.7\% using hand-crafted BVP
statistics under a different train/test split.
These results are therefore not directly comparable; our SVM serves as
an internal reference for the label-efficiency gap and a lower bound
for the supervised fine-tuning baselines.

Per-gesture accuracy averaged across folds (25\% labels) shows that
\textsc{Clap} is the most distinguishable class (77.94\%), while
\textsc{Draw-O-H} (44.29\%) and \textsc{Slide} (46.85\%) are the most
challenging, consistent with the fine-grained spatial trajectories of
those gestures.
We expect physics-informed representations to yield the greatest gains
on these two hard classes.

\subsubsection*{Widar\_LeNet baseline}
The Widar\_LeNet (CNN-5) architecture from SenseFi~\cite{ma2023sensefi},
adapted to our $T{=}20$ input, achieves $69.14 \pm 8.19\%$ with 25\%
labels and $79.15 \pm 7.43\%$ with 100\% labels — an improvement of
$+6.6$ and $+10.1$ percentage points over BVP+SVM respectively.
This confirms that convolutional feature extraction over the
spatial--velocity dimensions is substantially more informative than
flattened BVP vectors for cross-room generalisation.

A notable observation is that training accuracy reaches 100\% before
epoch~20 for all folds, yet test accuracy plateaus between 57\% and
75\%, indicating severe overfitting under the 25\%-label regime.
This overfitting gap motivates the use of unlabeled data via
self-supervised pre-training: PI-SSL leverages the full unlabeled pool
to learn a more general representation before supervised fine-tuning.

The 100\%-label result ($79.15\%$) remains $10.6$ points below the
Widar3.0 TPAMI reference (89.7\%), consistent with the architectural
differences between the two systems (CNN-5 vs.\ BVP-feature SVM in the
original paper).
The inter-fold variance ($\pm 8.19\%$) mirrors the pattern observed for
BVP+SVM, driven again by the Classroom fold's extreme train/test size
asymmetry.

\subsubsection*{Widar\_CNN-GRU baseline}
The Widar\_CNN-GRU architecture from SenseFi~\cite{ma2023sensefi},
adapted to $T{=}20$ with the final \texttt{Softmax} layer removed for
compatibility with \texttt{CrossEntropyLoss}, achieves $66.28 \pm 9.09\%$
with 25\% labels and $73.97 \pm 7.19\%$ with 100\% labels.

Contrary to expectation, CNN-GRU underperforms Widar\_LeNet by $2.9$
percentage points at 25\% labels, despite adding explicit temporal
modelling via the GRU layer.
Inspection of the training curves reveals the reason: CNN-GRU training
accuracy converges to $\sim\!88$--$91\%$ (never fully saturating), while
Widar\_LeNet reaches 100\% training accuracy before epoch~20.
The GRU introduces additional parameters that require more labeled
examples to tune, making the architecture less data-efficient in the
low-annotation regime.
At 100\% labels the gap narrows to $5.2$ points ($73.97\%$ vs
$79.15\%$), suggesting that CNN-GRU benefits more from additional data.
CNN-GRU also trains $11\times$ faster than Widar\_LeNet
(12.7~min vs.\ 144.8~min for all three folds), making it practical for
rapid experimentation.

\subsubsection*{Supervised baseline summary}
Table~\ref{tab:results} reveals three consistent patterns across all
supervised baselines:
(i)~every method suffers a $15$--$20$ percentage point accuracy drop on
Fold~1 (Classroom) relative to Folds~0 and~2, driven by the severe
train/test size asymmetry in that fold;
(ii)~increasing labeled data from 25\% to 100\% yields
$+6$--$10$ points for all architectures, indicating that label
efficiency is a genuine bottleneck;
(iii)~no supervised baseline approaches the Widar3.0 TPAMI reference
of 89.7\%, even with 100\% labels, confirming that the domain gap
between training and test rooms requires more than additional labeled
data alone.
These observations motivate PI-SSL's dual strategy: physics-informed
SSL pre-training to bridge the domain gap, and efficient fine-tuning to
exploit the limited labeled set.

\subsection{Ablation Study}
\label{sec:experiments:ablation}

[To be completed.]

% =============================================================================
\section{Conclusion}
\label{sec:conclusion}
% =============================================================================

[To be completed.]

% =============================================================================
% References
% =============================================================================

\begin{thebibliography}{9}

\bibitem{zheng2019widar3}
Y.~Zheng, Y.~Zhang, K.~Qian, G.~Zhang, Y.~Liu, C.~Wu, and Z.~Yang,
``Zero-Effort Cross-Domain Gesture Recognition with Wi-Fi,''
in \emph{Proc. 17th ACM Int. Conf. Mobile Systems, Applications, and
Services (MobiSys)}, Seoul, South Korea, Jun.\ 2019, pp.~313--325.

\bibitem{ma2023sensefi}
Y.~Ma, G.~Zhou, and S.~Wang,
``SenseFi: A Library and Benchmark on Deep-Learning-Empowered Wi-Fi
Human Sensing,''
\emph{Patterns}, vol.~4, no.~3, p.~100703, Mar.\ 2023.

\bibitem{yang2023arcfi}
H.~Yang \textit{et al.},
``ARC-Fi: Adaptive and Robust Cross-Environment Wi-Fi Sensing via
Semi-Supervised Contrastive Learning,''
\emph{arXiv preprint arXiv:2310.06328}, Oct.\ 2023.

\end{thebibliography}

\end{document}
